% GENERATED FILE. Edit canonical lessons, then run npm run generate:books.
% canonical-source-hash: fnv1a64:9dbec3f124c22ef8
% canonical-lessons: MR-C314-vacak, MR-C314-preksak, MR-C314-srota, MR-C314-sirsak, MR-C314-mukhprstha

\chapter{Readers and Viewers}
\label{ch:mr-readers-and-viewers}
\hlchaptermodality{314}

\begin{chapteropening}
\textbf{By the end of this chapter:} \emph{I can talk about the people who read, watch and listen.}

Complete the last lesson of chapter 314: I can talk about the people who read, watch and listen.
\end{chapteropening}

\section[\texorpdfstring{\mr{वाचक} (vācak)}{वाचक (vācak)}]{\mr{वाचक} (vācak) — a reader}

\label{lesson:MR-C314-vacak}



\begin{quote}
Before the new one: say the Marathi for prominent, then the Marathi for to proclaim.
\end{quote}

\subsection*{You'll want to know: \mr{वाचक}}
\textbf{\mr{वाचक}} — \emph{vācak} — \textquotedblleft{}a reader\textquotedblright{}.

\begin{quote}
\mr{या} \mr{मासिकाचे} \mr{हजारो} \mr{वाचक} \mr{आहेत}.
\end{quote}

\begin{quote}
\emph{yā māsikāce hajāro vācak āhet.}
\end{quote}

\begin{quote}
This magazine has thousands of readers.
\end{quote}

Masculine.

Plural \mr{वाचक}.

\begin{cousinweb}
\textbf{\mr{वाचक}} (\emph{vācak}) is a tatsama, taken from Sanskrit \emph{vācaka}, 'speaking, expressing' (Monier-Williams); Marathi uses it for the one who reads, from \emph{vācṇe}, 'to read'.
\end{cousinweb}

\begin{grammarlens}[title={What you've built}]
One more word for telling what happened.
\end{grammarlens}

\subsection*{Your turn}
\begin{itemize}
\raggedright
  \item \emph{Say it:} \emph{vācak}
  \item \emph{Say it:} \emph{vācak}, once more
  \item \emph{Say it:} \emph{yā māsikāce hajāro vācak āhet.} — This magazine has thousands of readers.
  \item \emph{Recall:} say the Marathi for prominent, then the Marathi for to proclaim, then say \emph{vācak} again
\end{itemize}

\begin{tcolorbox}[breakable,colback=teal!4,colframe=teal!35!black,title={Before you move on}]
What does \mr{वाचक} mean? (\textquotedblleft{}A reader\textquotedblright{}.) Say it once more.
\end{tcolorbox}
\section[\texorpdfstring{\mr{प्रेक्षक} (prekṣak)}{प्रेक्षक (prekṣak)}]{\mr{प्रेक्षक} (prekṣak) — a viewer; a spectator}

\label{lesson:MR-C314-preksak}



\begin{quote}
Before the new one: say the Marathi for to proclaim, then the Marathi for a reader.
\end{quote}

\subsection*{You'll want to know: \mr{प्रेक्षक}}
\textbf{\mr{प्रेक्षक}} — \emph{prekṣak} — \textquotedblleft{}a viewer; a spectator\textquotedblright{}.

\begin{quote}
\mr{नाटक} \mr{पाहायला} \mr{हजारो} \mr{प्रेक्षक} \mr{आले} \mr{होते}.
\end{quote}

\begin{quote}
\emph{nāṭak pāhāylā hajāro prekṣak āle hote.}
\end{quote}

\begin{quote}
Thousands of spectators had come to watch the play.
\end{quote}

Masculine.

Plural \mr{प्रेक्षक}.

\begin{cousinweb}
\textbf{\mr{प्रेक्षक}} (\emph{prekṣak}) is a tatsama, taken from Sanskrit \emph{prekṣaka}, 'a looker-on, a spectator' (Monier-Williams).
\end{cousinweb}

\begin{grammarlens}[title={What you've built}]
One more word for telling what happened.
\end{grammarlens}

\subsection*{Your turn}
\begin{itemize}
\raggedright
  \item \emph{Say it:} \emph{prekṣak}
  \item \emph{Say it:} \emph{prekṣak}, once more
  \item \emph{Say it:} \emph{nāṭak pāhāylā hajāro prekṣak āle hote.} — Thousands of spectators had come to watch the play.
  \item \emph{Recall:} say the Marathi for to proclaim, then the Marathi for a reader, then say \emph{prekṣak} again
\end{itemize}

\begin{tcolorbox}[breakable,colback=teal!4,colframe=teal!35!black,title={Before you move on}]
What does \mr{प्रेक्षक} mean? (\textquotedblleft{}A viewer; a spectator\textquotedblright{}.) Say it once more.
\end{tcolorbox}
\section[\texorpdfstring{\mr{श्रोता} (śrotā)}{श्रोता (śrotā)}]{\mr{श्रोता} (śrotā) — a listener}

\label{lesson:MR-C314-srota}



\begin{quote}
Before the new one: say the Marathi for a reader, then the Marathi for a viewer.
\end{quote}

\subsection*{You'll want to know: \mr{श्रोता}}
\textbf{\mr{श्रोता}} — \emph{śrotā} — \textquotedblleft{}a listener\textquotedblright{}.

\begin{quote}
\mr{कार्यक्रमानंतर} \mr{एका} \mr{श्रोत्याने} \mr{प्रश्न} \mr{विचारला}.
\end{quote}

\begin{quote}
\emph{kāryakramānantar ekā śrotyāne praśna vicārlā.}
\end{quote}

\begin{quote}
After the programme a listener asked a question.
\end{quote}

Masculine. Plural: \mr{श्रोते}, \emph{śrote}.

The oblique is \mr{श्रोत्या}-, the plural \mr{श्रोते}.

\begin{cousinweb}
\textbf{\mr{श्रोता}} (\emph{śrotā}) is a tatsama, taken from Sanskrit \emph{śrotṛ}, 'one who hears, a hearer' (Monier-Williams), in its nominative form \emph{śrotā}.
\end{cousinweb}

\begin{grammarlens}[title={What you've built}]
One more word for telling what happened.
\end{grammarlens}

\subsection*{Your turn}
\begin{itemize}
\raggedright
  \item \emph{Say it:} \emph{śrotā}
  \item \emph{Say it:} \emph{śrotā}, once more
  \item \emph{Say it:} \emph{kāryakramānantar ekā śrotyāne praśna vicārlā.} — After the programme a listener asked a question.
  \item \emph{Recall:} say the Marathi for a reader, then the Marathi for a viewer, then say \emph{śrotā} again
\end{itemize}

\begin{tcolorbox}[breakable,colback=teal!4,colframe=teal!35!black,title={Before you move on}]
What does \mr{श्रोता} mean? (\textquotedblleft{}A listener\textquotedblright{}.) Say it once more.
\end{tcolorbox}
\section[\texorpdfstring{\mr{शीर्षक} (śīrṣak)}{शीर्षक (śīrṣak)}]{\mr{शीर्षक} (śīrṣak) — a title; a heading}

\label{lesson:MR-C314-sirsak}



\begin{quote}
Before the new one: say the Marathi for a viewer, then the Marathi for a listener.
\end{quote}

\subsection*{You'll want to know: \mr{शीर्षक}}
\textbf{\mr{शीर्षक}} — \emph{śīrṣak} — \textquotedblleft{}a title; a heading\textquotedblright{}.

\begin{quote}
\mr{या} \mr{कवितेचं} \mr{शीर्षक} \mr{फार} \mr{सुंदर} \mr{आहे}.
\end{quote}

\begin{quote}
\emph{yā kavitecaṁ śīrṣak phār sundar āhe.}
\end{quote}

\begin{quote}
The title of this poem is very beautiful.
\end{quote}

Neuter.

\mr{शीर्षकाखाली}, under the heading.

\begin{cousinweb}
\textbf{\mr{शीर्षक}} (\emph{śīrṣak}) is built on Sanskrit \emph{śīrṣa}, 'the head' (Monier-Williams): what stands at the head of a page.
\end{cousinweb}

\begin{grammarlens}[title={What you've built}]
One more word for telling what happened.
\end{grammarlens}

\subsection*{Your turn}
\begin{itemize}
\raggedright
  \item \emph{Say it:} \emph{śīrṣak}
  \item \emph{Say it:} \emph{śīrṣak}, once more
  \item \emph{Say it:} \emph{yā kavitecaṁ śīrṣak phār sundar āhe.} — The title of this poem is very beautiful.
  \item \emph{Recall:} say the Marathi for a viewer, then the Marathi for a listener, then say \emph{śīrṣak} again
\end{itemize}

\begin{tcolorbox}[breakable,colback=teal!4,colframe=teal!35!black,title={Before you move on}]
What does \mr{शीर्षक} mean? (\textquotedblleft{}A title; a heading\textquotedblright{}.) Say it once more.
\end{tcolorbox}
\section[\texorpdfstring{\mr{मुखपृष्ठ} (mukhpṛṣṭha)}{मुखपृष्ठ (mukhpṛṣṭha)}]{\mr{मुखपृष्ठ} (mukhpṛṣṭha) — a front page; a cover}

\label{lesson:MR-C314-mukhprstha}



\begin{quote}
Before the new one: say the Marathi for a listener, then the Marathi for a title.
\end{quote}

\subsection*{You'll want to know: \mr{मुखपृष्ठ}}
\textbf{\mr{मुखपृष्ठ}} — \emph{mukhpṛṣṭha} — \textquotedblleft{}a front page; a cover\textquotedblright{}.

\begin{quote}
\mr{पुस्तकाच्या} \mr{मुखपृष्ठावर} \mr{एक} \mr{जुना} \mr{किल्ला} \mr{आहे}.
\end{quote}

\begin{quote}
\emph{pustakācyā mukhpṛṣṭhāvar ek junā killā āhe.}
\end{quote}

\begin{quote}
There is an old fort on the book's cover.
\end{quote}

Neuter.

Built from \mr{मुख}, face, and \mr{पृष्ठ}, a page.

\begin{cousinweb}
\textbf{\mr{मुखपृष्ठ}} (\emph{mukhpṛṣṭha}) is a modern compound built from Sanskrit \emph{mukha}, 'the face', and \emph{pṛṣṭha}, 'the back, a surface, a page' (Monier-Williams).
\end{cousinweb}

\begin{grammarlens}[title={What you've built}]
One more word for telling what happened.
\end{grammarlens}

\subsection*{Your turn}
\begin{itemize}
\raggedright
  \item \emph{Say it:} \emph{mukhpṛṣṭha}
  \item \emph{Say it:} \emph{mukhpṛṣṭha}, once more
  \item \emph{Say it:} \emph{pustakācyā mukhpṛṣṭhāvar ek junā killā āhe.} — There is an old fort on the book's cover.
  \item \emph{Recall:} say the Marathi for a listener, then the Marathi for a title, then say \emph{mukhpṛṣṭha} again
\end{itemize}

\begin{tcolorbox}[breakable,colback=teal!4,colframe=teal!35!black,title={Before you move on}]
What does \mr{मुखपृष्ठ} mean? (\textquotedblleft{}A front page; a cover\textquotedblright{}.) Say it once more.
\end{tcolorbox}
